% =====================================================================
%  CHƯƠNG 7 — KẾT LUẬN VÀ HƯỚNG PHÁT TRIỂN
% =====================================================================
\chapter{Kết luận và hướng phát triển}\label{chap:ketluan}

\section{Kết luận}
Tiểu luận đã nghiên cứu đồ thị tri thức, sinh tăng cường truy xuất và tác tử thông minh, rồi hiện thực hóa toàn bộ
thành một hệ thống chạy được cho hai bài toán của Việt Nam: tư vấn cho người bệnh đái tháo đường típ 2 và phân tích
tác động vĩ mô lên cổ phiếu VN30 và vàng SJC. Mọi con số trong báo cáo được sinh tự động từ kết quả thí nghiệm, và mọi
đoạn mã được trích từ chính mã nguồn đã chạy ra các con số đó. Kết quả chính:

\begin{enumerate}
  \item \textbf{Đồ thị tri thức.} Hai KG có bản thể học riêng, ràng buộc miền -- miền giá trị và nguồn văn bản cho
        từng cạnh: \texttt{VietnamDiabetesKG} (\SL{kg/health/nodes} thực thể, \SL{kg/health/edges} quan hệ) và
        \texttt{VietnamFinanceKG} (\SL{kg/finance/nodes} thực thể, \SL{kg/finance/edges} quan hệ). Suy diễn RDFS làm
        \SL{dem/sparql} truy vấn SPARQL trả đúng kết quả, trong đó có truy vấn chỉ có kết quả \emph{sau} suy diễn. Với
        câu hỏi có điều kiện số trên dữ liệu bảng, truy vấn KG trả đúng
        \SL{kg/tabular/numeric_query/kg_hits}/\SL{kg/tabular/numeric_query/truth} hồ sơ, trong khi tìm kiếm văn bản
        TF-IDF chỉ đưa \SL{kg/tabular/numeric_query/tfidf_top20_correct} hồ sơ đúng vào top-20.
  \item \textbf{Truy xuất.} Trên \SL{retrieval_grid/n_questions} câu hỏi, chia đoạn theo cấu trúc kết hợp BM25,
        e5 và xếp hạng lại đạt Hit@1 \SL{retrieval_grid/grid/structure/hybrid_rerank/hit@1Pct}\%, so với
        \SL{retrieval_grid/grid/fixed/dense/hit@1Pct}\% của chia đoạn cố định với dense thuần. Hai kỹ thuật RAG nâng
        cao phổ biến lại \emph{làm giảm} chất lượng khi dùng LLM 1,5B: HyDE hạ Hit@1 từ
        \SL{advanced/agg/dense/hit@1Pct}\% xuống \SL{advanced/agg/hyde/hit@1Pct}\%, Multi-Query xuống
        \SL{advanced/agg/multiquery/hit@1Pct}\%. GraphRAG cục bộ nâng độ phủ đáp án từ
        \SL{graphrag/agg/naive/ans_covPct}\% lên \SL{graphrag/agg/graph_local/ans_covPct}\%.
  \item \textbf{Học sâu.} MLP nguy cơ đái tháo đường đạt ROC-AUC \SL{diabetes/mlp/roc_aucPct}\% và PR-AUC
        \SL{diabetes/mlp/pr_aucPct}\% trên tập kiểm thử. BiLSTM dự báo lợi suất không thắng có ý nghĩa dự báo ngây thơ
        (\SL{dem/fc_beat}/\SL{dem/fc_models} mô hình có tỉ số RMSE dưới 1, tốt nhất chỉ giảm
        \SL{dem/fc_best_gainPct}\% sai số), phù hợp giả thuyết thị trường hiệu quả dạng yếu.
  \item \textbf{Tác tử.} Tác tử ReAct với \SL{dem/tools} công cụ gọi đủ công cụ cần thiết ở mọi kịch bản và nâng độ
        phủ dữ kiện từ \SL{agent/agg/all/coverage_llm_onlyPct}\% (LLM trả lời một mình) lên
        \SL{agent/agg/all/coveragePct}\%, đồng thời loại bỏ các lỗi nguy hiểm của LLM như khẳng định metformin không
        ảnh hưởng khi uống bia hay đề xuất thuốc chống chỉ định cho người suy tim.
  \item \textbf{Ứng dụng.} Web App Flask chạy toàn bộ mô hình cục bộ trên một GPU 8 GB, hiển thị vết suy luận, câu
        trả lời có trích dẫn, đồ thị tri thức tương tác và so sánh ba bộ truy xuất.
\end{enumerate}

Bài học chung của tiểu luận là các thành phần bổ trợ nhau đúng như lý thuyết dự đoán: KG cho ràng buộc chính xác và
suy luận nhiều bước, RAG cho văn bản gốc để giải thích và trích dẫn, mô hình học sâu cho đánh giá định lượng từ dữ
liệu của từng người, còn LLM chỉ làm việc nó làm tốt nhất là diễn đạt. Khi LLM được giao thêm việc nhớ tri thức, tự
viết truy vấn giả định (HyDE) hay tự chọn công cụ, chất lượng giảm rõ rệt với mô hình nhỏ.

\section{Hạn chế}
\begin{itemize}
  \item \textbf{LLM nhỏ không tuân thủ định dạng trích dẫn.} Chỉ \SL{agent/agg/all/llm_passPct}\% bản nháp qua
        kiểm chứng; \SL{dem/agent_nocite} bản bị từ chối vì không có số hiệu [n]. Câu trả lời cuối vì vậy thường là
        bản tổng hợp theo khuôn: đúng và truy vết được, nhưng kém tự nhiên.
  \item \textbf{Bộ kiểm chứng chỉ kiểm con số và trích dẫn.} Nó không phát hiện được câu trộn cơ chế hay viện dẫn
        ngoài bằng chứng khi câu đó không chứa con số (Mục~\ref{sec:agent-ketqua}).
  \item \textbf{Xác suất của MLP chưa hiệu chỉnh.} Trọng số lớp dương làm xác suất bị đẩy cao so với tần suất thực;
        con số ``xác suất mắc bệnh'' chỉ nên đọc như điểm xếp hạng nguy cơ.
  \item \textbf{KG viết tay và nhỏ.} Hai KG có vài chục thực thể, cạnh được viết từ kho văn bản thay vì trích xuất
        tự động; dấu cạnh tài chính cố định trong khi tương quan thực tế đổi dấu theo giai đoạn, như cạnh USD/VND $\to$
        FPT chỉ khớp dữ liệu khi xét độ trễ.
  \item \textbf{Bộ lập kế hoạch dựa trên luật.} Nó tái lập được và chính xác trên tập kịch bản, nhưng không thích
        nghi với câu hỏi ngoài mẫu, và xét vai trò nguyên nhân -- kết quả theo lớp của thực thể thay vì theo câu.
  \item \textbf{Quy mô đánh giá.} \SL{retrieval_grid/n_questions} câu hỏi truy xuất, \SL{dem/grag_q} câu hỏi GraphRAG
        và \SL{dem/agent_scen} kịch bản tác tử đủ để so sánh các phương án trong bài, nhưng chưa đủ để kết luận có ý
        nghĩa thống kê cho những chênh lệch nhỏ.
\end{itemize}

\section{Hướng phát triển}
\begin{itemize}
  \item \textbf{LLM tuân thủ trích dẫn tốt hơn}: dùng mô hình 7B lượng tử hóa 4-bit vừa với GPU 8 GB, hoặc tinh chỉnh
        LoRA mô hình 1,5B trên các cặp (bằng chứng, câu trả lời có trích dẫn) sinh từ chính bản tổng hợp theo khuôn.
  \item \textbf{Kiểm chứng theo ngữ nghĩa}: đối chiếu từng câu của câu trả lời với bằng chứng được trích bằng một mô
        hình suy diễn ngôn ngữ tự nhiên (NLI) đa ngữ, thay cho chỉ kiểm con số.
  \item \textbf{Hiệu chỉnh xác suất}: áp dụng temperature scaling hoặc hồi quy isotonic trên tập kiểm định để xác
        suất của MLP phản ánh đúng tần suất mắc bệnh.
  \item \textbf{Nối tầng học sâu với KG}: đối chiếu chỉ số xét nghiệm của hồ sơ với các ngưỡng chẩn đoán đã có trong
        KG, để tác tử tự kết luận ``tiền đái tháo đường'' thay vì chỉ báo điểm nguy cơ.
  \item \textbf{KG tài chính có bằng chứng thực nghiệm}: gắn tương quan, độ trễ và khoảng thời gian hiệu lực vào thuộc
        tính của từng cạnh, cập nhật định kỳ từ dữ liệu, để tác tử nói ``tác động cùng chiều thể hiện sau khoảng hai
        tháng'' thay vì một dấu cố định.
  \item \textbf{Mở rộng KG tự động}: dùng LLM trích xuất bộ ba từ văn bản mới của Bộ Y tế, NHNN và báo cáo doanh
        nghiệp, kiểm tra bằng ràng buộc bản thể học trước khi nạp, và chuyển lưu trữ từ NetworkX sang cơ sở dữ liệu đồ
        thị như Neo4j khi quy mô tăng.
  \item \textbf{Bộ lập kế hoạch lai}: giữ luật làm mặc định, nhưng cho LLM đề xuất thêm công cụ khi luật không khớp,
        với bộ kiểm tra tên và tham số trước khi thực thi.
\end{itemize}
